\begin{helper}
Let \(I\subseteq[n]\) have \(|I|=k\), with
\[
k=\noderef{TwoBiteNaturalIndependenceScale}(1+\varepsilon,n),
\]
and assume \(0\le\varepsilon_1\),
\(\noderef{ParameterHierarchyEventualComparisons}
(\varepsilon,\varepsilon_1,\varepsilon_2,n_0)\), \(n>n_0\),
\[
m=\noderef{TwoBiteNaturalReducedVertexCount}(n),\qquad
p=\frac12\sqrt{\frac{\log n}{n}},
\]
and
\(\noderef{TwoBiteRegularityEvent}(k,\varepsilon,\varepsilon_1,\varepsilon_2,
\omega)\).  Put
\[
T=\noderef{TwoBiteStagedOpenPairs}(\omega,\varepsilon,I).
\]
For a tagged coordinate \(e\), let \(e\) be present mean that in the red case
the corresponding edge is present in \(\noderef{TwoBiteRedGraph}(\omega)\),
and in the blue case it is present in
\(\noderef{TwoBiteBlueGraph}(\omega)\).

Define a red-exceptional predicate as follows.  A coordinate
\(\operatorname{red}(q)\) is red-exceptional exactly when there is a blue base
vertex \(b\) which is large, medium, or small for \(I\), and
\[
q\in\noderef{TwoBiteBasePairs}
\bigl(\pi_R^\omega(\noderef{TwoBiteX}(\omega,I,\operatorname{blue}(b)))\bigr);
\]
blue-tagged coordinates are not red-exceptional.  Define the
blue-exceptional predicate symmetrically: a coordinate
\(\operatorname{blue}(q)\) is blue-exceptional exactly when there is a red base
vertex \(r\) which is large, medium, or small for \(I\), and
\[
q\in\noderef{TwoBiteBasePairs}
\bigl(\pi_B^\omega(\noderef{TwoBiteX}(\omega,I,\operatorname{red}(r)))\bigr);
\]
red-tagged coordinates are not blue-exceptional.

Then there exist finite sets \(E_R,E_B\) of tagged projected coordinates such
that
\[
e\in E_R\quad\Longleftrightarrow\quad e\in T
\text{ and }e\text{ is red-exceptional},
\]
\[
e\in E_B\quad\Longleftrightarrow\quad e\in T
\text{ and }e\text{ is blue-exceptional}.
\]
If \(I\) is independent in the shadow graph of
\(\noderef{TwoBiteFinalGraph}(\omega)\), then every present
staged-open coordinate \(e\in T\) is red-exceptional or blue-exceptional.
Moreover
\[
\noderef{ClosedPairsBound}(|E_R|,3\varepsilon_1,k)
\qquad\text{and}\qquad
\noderef{ClosedPairsBound}(|E_B|,3\varepsilon_1,k).
\]
\end{helper}
\begin{proof}
Define
\[
E_R=\{e\in T:e\hbox{ is red-exceptional}\},
\qquad
E_B=\{e\in T:e\hbox{ is blue-exceptional}\}.
\]
These are finite filters of the finite set \(T\), so they are legitimate
finite sets.  Their defining membership equivalences are exactly the first two
conjuncts in the statement.

We next prove the classification of present staged-open coordinates.  We give
the red case in detail; the blue case is obtained by interchanging the two
colors.  Let \(e=\operatorname{red}(q)\in T\), suppose that the red base edge
\(q\) is present, and assume that \(I\) is independent in the final graph.
Unfolding \(\noderef{TwoBiteStagedOpenPairs}\) gives
\[
e\in\noderef{TwoBiteProjectionPairSet}(\omega,I),\qquad
\neg\noderef{TwoBiteProjectionPairSameColorClosed}(\omega,I,e),\qquad
\neg\noderef{TwoBiteProjectionPairTouched}(\omega,\varepsilon,I,e).
\tag{1}
\]
The projected-pair membership supplies final vertices \(u,v\in I\) whose red
projections form the increasing base pair \(q\), in the sense of
\(\noderef{TwoBiteBasePairs}\).  Since \(q\) is present, the overlay contains
the red edge \(uv\).  Independence of \(I\) in
\(\noderef{TwoBiteFinalGraph}(\omega)\) says that this edge does not remain in
the final graph.  Therefore, by the deletion rules in
\(\noderef{TwoBiteFinalGraph}\), the present red edge \(uv\) is selected by
one of the red deletion alternatives.

First consider a same-color red deletion.  If \(q\) is the latest red
coordinate of the deleting red triangle, then the third vertex of the
triangle is a red base vertex \(r\) with
\[
u,v\in\noderef{TwoBiteXPlus}(\omega,I,\operatorname{red}(r)).
\]
Equivalently,
\[
q\in\noderef{TwoBiteBasePairs}
\bigl(\pi_R^\omega(\noderef{TwoBiteXPlus}
(\omega,I,\operatorname{red}(r)))\bigr),
\]
which is precisely
\(\noderef{TwoBiteProjectionPairSameColorClosed}(\omega,I,e)\), contrary to
(1).  If the deleting red triangle has a red coordinate later than \(q\), the
deletion rule removes that later coordinate, not \(q\), so it cannot explain
the absence of the present edge \(uv\).  If the deletion has already forced
one endpoint of \(q\) to be staged-revealed before the current coordinate is
tested, then the touched boundary clause of
\(\noderef{FixedSetClosedPredicateBoundary}\) gives
\(\noderef{TwoBiteProjectionPairTouched}(\omega,\varepsilon,I,e)\), again
contrary to (1).  Thus no same-color or already-touched deletion alternative
can be responsible for this present staged-open red coordinate.

The remaining possible deletion of a red edge is a mixed triangle with a blue
witness.  Let \(b\) be the blue base vertex supplied by that mixed triangle.
Both \(u\) and \(v\) are blue-neighbors of \(b\), hence
\[
u,v\in\noderef{TwoBiteX}(\omega,I,\operatorname{blue}(b)).
\]
Projecting these two vertices to the red copy gives
\[
q\in\noderef{TwoBiteBasePairs}
\bigl(\pi_R^\omega(\noderef{TwoBiteX}
(\omega,I,\operatorname{blue}(b)))\bigr). \tag{2}
\]
The set \(\noderef{TwoBiteX}(\omega,I,\operatorname{blue}(b))\) is a subset
of \(I\), so its cardinality lies between \(0\) and \(|I|=k\).  Comparing this
cardinality with the small, medium, large, and huge cutoffs gives exactly one
of the four alternatives
\[
b\in S_I,\qquad b\in M_I,\qquad b\in L_I,\qquad b\in H_I.
\tag{3}
\]
If \(b\in H_I\), then (2) is a huge opposite-color witness for the red
coordinate.  The regularity event supplies the fiber-degree event, and the
statement already carries the eventual-comparisons hypothesis, \(n>n_0\), and
the rounded \(m,p\) identities and the scale identity
\(|I|=k=\noderef{TwoBiteNaturalIndependenceScale}(1+\varepsilon,n)\).
Therefore
\(\noderef{FixedSetHugeOppositeWitnessTouched}\) applies to this witness.  It
says that \(\operatorname{red}(q)=e\) is pre-terminal recorded, and hence is
not in
\(\noderef{TwoBiteStagedOpenPairs}(\omega,\varepsilon,I)=T\).  This
contradicts the standing assumption \(e\in T\).  Hence \(b\) is not huge.
From (3), \(b\) is large, medium, or small, and (2) is exactly the
red-exceptional witness required in the statement.  Therefore \(e\) is
red-exceptional.

The blue case repeats the same argument with the two colors reversed.  A
present staged-open coordinate \(e=\operatorname{blue}(q)\) cannot be deleted
by a latest same-color blue triangle, because that would be
\(\noderef{TwoBiteProjectionPairSameColorClosed}\), and it cannot be deleted
through a previously revealed endpoint, because that would be
\(\noderef{TwoBiteProjectionPairTouched}\).  Thus it has a mixed red witness
\(r\) with
\[
q\in\noderef{TwoBiteBasePairs}
\bigl(\pi_B^\omega(\noderef{TwoBiteX}
(\omega,I,\operatorname{red}(r)))\bigr).
\]
The cutoff partition puts \(r\) in small, medium, large, or huge; the huge
case is excluded by the blue half of
\(\noderef{FixedSetHugeOppositeWitnessTouched}\).  With the same
regularity-derived fiber-degree event and the same hierarchy, rounded
parameter, and scale hypotheses, that helper says that a blue coordinate with
such a relevant huge red witness is pre-terminal recorded and therefore is
not in
\(\noderef{TwoBiteStagedOpenPairs}(\omega,\varepsilon,I)=T\).  Hence \(r\)
is large, medium, or small, and the coordinate is blue-exceptional.  This
proves that every present staged-open coordinate needed by final independence
belongs to \(E_R\cup E_B\), with the tag deciding which exceptional predicate
can hold.

It remains to count the two exceptional sets.  We count \(E_R\); the blue
count is symmetric.  For each \(e=\operatorname{red}(q)\in E_R\), choose one
blue witness \(b(e)\) and one class label
\[
\mathcal A(e)\in\{L_I,M_I,S_I\}
\]
certifying that \(b(e)\) is large, medium, or small and that
\[
q\in\noderef{TwoBiteBasePairs}
\bigl(\pi_R^\omega(\noderef{TwoBiteX}
(\omega,I,\operatorname{blue}(b(e))))\bigr). \tag{4}
\]
If several witnesses or labels exist, choose one by finite choice.  Let
\(E_R^L,E_R^M,E_R^S\) be the three subfamilies selected by this label.  These
subfamilies are disjoint and cover \(E_R\).

Fix one class \(\mathcal A\in\{L_I,M_I,S_I\}\) and one
\(e=\operatorname{red}(q)\in E_R^\mathcal A\).  From (4) choose final vertices
\(u,v\in\noderef{TwoBiteX}(\omega,I,\operatorname{blue}(b(e)))\) whose red
projected base pair is \(q\).  Since the chosen witness \(b(e)\) lies in
\(\mathcal A\), the unordered final pair \(\{u,v\}\) belongs to
\[
\noderef{TwoBiteClosedPairsInClass}(\omega,I,\mathcal A).
\tag{5}
\]
The map \(e\mapsto\{u,v\}\) has multiplicity at most one.  Indeed, once a
final unordered pair \(\{u,v\}\) is fixed, the red projected coordinate
generated by that pair is uniquely the increasing representative
\(\{\pi_R^\omega(u),\pi_R^\omega(v)\}\) selected by
\(\noderef{TwoBiteBasePairs}\).  Therefore two red tagged coordinates mapped
to the same final pair must be the same tagged coordinate.  Hence
\[
|E_R^\mathcal A|
\le
|\noderef{TwoBiteClosedPairsInClass}(\omega,I,\mathcal A)|.
\tag{6}
\]

Apply (6) to the three classes.  For \(\mathcal A=L_I\),
\(\noderef{LargeClosedPairsBound}\), with the present eventual-comparisons
hypothesis, \(n>n_0\), the scale identity
\[
|I|=k=\noderef{TwoBiteNaturalIndependenceScale}(1+\varepsilon,n),
\]
and the fiber-degree conjunct of
\(\noderef{TwoBiteRegularityEvent}\), gives
\[
\noderef{ClosedPairsBound}
\bigl(|\noderef{TwoBiteClosedPairsInClass}(\omega,I,L_I)|,\varepsilon_1,k\bigr).
\tag{7}
\]
The medium and small conjuncts of
\(\noderef{TwoBiteRegularityEvent}\), namely
\(\noderef{TwoBiteMediumClosedPairsEvent}\) and
\(\noderef{TwoBiteSmallClosedPairsEvent}\), give the analogous bounds for
\(M_I\) and \(S_I\):
\[
\noderef{ClosedPairsBound}
\bigl(|\noderef{TwoBiteClosedPairsInClass}(\omega,I,M_I)|,\varepsilon_1,k\bigr),
\qquad
\noderef{ClosedPairsBound}
\bigl(|\noderef{TwoBiteClosedPairsInClass}(\omega,I,S_I)|,\varepsilon_1,k\bigr).
\tag{8}
\]
Unfolding \(\noderef{ClosedPairsBound}\), (7) and (8) say that each of the
three right-hand sides in (6) is at most \(\varepsilon_1k^2\).  Since
\(E_R\) is the disjoint union of \(E_R^L,E_R^M,E_R^S\),
\[
|E_R|
\le \varepsilon_1k^2+\varepsilon_1k^2+\varepsilon_1k^2
=3\varepsilon_1k^2.
\]
This is exactly
\(\noderef{ClosedPairsBound}(|E_R|,3\varepsilon_1,k)\).  Repeating the same
construction with red and blue interchanged maps \(E_B\) into closed pairs in
the large, medium, and small red-witness classes, with the same multiplicity
one property from \(\noderef{TwoBiteBasePairs}\).  The same three regularity
and large-class bounds give
\[
|E_B|\le3\varepsilon_1k^2,
\]
which is the asserted
\(\noderef{ClosedPairsBound}(|E_B|,3\varepsilon_1,k)\).
\end{proof}
